\handbooktitle{EF Core plugin development · reference and practice}{The EF Core Plugin Handbook}
  {From the options extension to the NuGet package: everything you need for professional and
  even very complex extensions of Entity Framework Core. Explained throughout with a
  complete multi-tenancy plugin for PostgreSQL and SQL Server.}
\begin{handbookmeta}
  Edition & \handbookversion \\
  Versions & EF Core 10, Npgsql.EntityFrameworkCore.PostgreSQL 10, .NET 10 \\
  Sample & Complete solution at \textbf{github.com/guenni81/EF-Core-Plugin-Handbook}: core
  package, two provider packages, two sample apps, unit and integration tests \\
  Verified & The sample code is built against EF Core 10 and Npgsql 10; unit and integration
  tests run against real PostgreSQL and SQL Server. \\
  Structure & 9 parts, 42 chapters, 4 appendices \\
  License & MIT License, © 2026 guenni81. Free to use, copy, modify and share, including
  commercially; see the last page. \\
\end{handbookmeta}

\section*{About this book}

Many EF Core extensions start out as a handful of extension methods and then grow without a
plan. If you understand how EF Core is built internally, however, you can write extensions
that feel like part of EF Core itself: they are configured through the options, run inside
the internal service provider, translate custom functions to SQL, generate migrations, write
readable snapshots, and combine cleanly with other plugins.

This book explains each of these building blocks in two layers. First comes the
\textbf{principle}: how does the mechanism work in EF Core, which variants exist, and what
do you need to watch out for? Then comes the \textbf{implementation} in the sample plugin,
with the code that is also in the repository. The chapters build on each other, but can also
be read individually as a reference. The code in the repository and the snippets in this
edition use English comments and messages. Each snippet's label names the file.

\begin{datatable}{H{0.50}H{1.50}}
  \textbf{Box} & \textbf{Meaning} \\
  \tablehead
  \textcolor[HTML]{D99A00}{\textbf{Important}} & Pitfalls that lead to incorrect behavior or
  security holes. \\
  \tablerule
  \textcolor[HTML]{2E8B57}{\textbf{Tip}} & Proven practice that saves work or improves
  quality. \\
  \tablerule
  \textcolor[HTML]{1D4E89}{\textbf{Note}} & Background, context and references to other
  chapters. \\
\end{datatable}

\begin{note}[Public and internal APIs]
  Some APIs that plugins need live in namespaces ending in \code{.\allowbreak{}Internal}. EF
  Core calls them “pubternal”: publicly visible, but without compatibility guarantees. The
  book marks these places, and chapter~\ref{ch:versioning-and-compatibility} explains how to
  deal with them.
\end{note}

\tableofcontents

% ============================================================
\setpartpreamble{\partintro
  {What a plugin is, how EF Core is built internally, and which scenario accompanies us
  through the book.}
  {
    \item Chapter~\ref{ch:what-is-an-ef-core}: What is an EF Core plugin?
    \item Chapter~\ref{ch:ef-core-s-architecture-from}: EF Core's architecture from a
    plugin's point of view
    \item Chapter~\ref{ch:the-sample-scenario}: The sample scenario
  }}
\part{Foundations}\label{part:I}

% ------------------------------------------------------------
\chapter{What is an EF Core plugin?}\label{ch:what-is-an-ef-core}

The term “plugin” is used for very different things. For this book, it helps to distinguish
four levels that differ in how deeply they reach into EF Core:

\begin{datatable}{H{0.65}H{1.41}H{0.94}}
  \textbf{Level} & \textbf{Description} & \textbf{Examples} \\
  \tablehead
  1 · Helpers & Extension methods on \code{Db\allowbreak{}Context} or \code{IQueryable} that
  only call public API. & Paging helpers, \code{Where\allowbreak{}If}, seeding methods \\
  \tablerule
  2 · Model helpers & Fluent API extensions and annotations, without services of their own
  inside EF Core. & Custom \code{Has\allowbreak{}Xyz(\allowbreak{})} methods, ModelBuilder
  helpers \\
  \tablerule
  3 · Real plugins & Own services in the internal service provider: conventions, query
  translation, type mappings, migrations, design time. & NetTopologySuite plugin, NodaTime
  plugin, EFCore.NamingConventions \\
  \tablerule
  4 · Providers & A complete database integration including connection, update pipeline and
  SQL generation. & Npgsql, Pomelo (MySQL), SQLite, Cosmos \\
\end{datatable}

This book covers \textbf{level 3} completely. Levels 1 and 2 need no book of their own: they
only use EF's public API, and chapter~\ref{ch:annotations-and-fluent-api} covers what is
specific to Fluent API extensions. The tools of level 4 (such as replacing the SQL
generator) are explained where plugins need them.

\section{Who this book is for}

For developers who already use EF Core confidently: \code{Db\allowbreak{}Context},
migrations, LINQ queries and dependency injection are assumed. No knowledge of EF's
internals is required; the book introduces them where they are needed. Experience with
PostgreSQL helps with the sample, but every SQL construct is explained when it first
appears. Chapter~\ref{ch:ef-core-s-architecture-from} ends with a complete mini plugin that
shows all building blocks on a single page before the large sample begins.

\section{What a level 3 plugin can do}

A complete plugin can extend EF Core at practically every point: it can make new CLR types
available as column types, translate custom functions in LINQ to SQL, adapt the model
automatically through conventions, accept metadata through the Fluent API and attributes,
validate the model, intervene when saving and when opening connections, log its own events,
add custom SQL to migrations and extend the \code{dotnet ef} tooling. The sample in this
book uses every one of these possibilities.

\begin{tip}[Role models in the ecosystem]
  Before writing your own, look at existing plugins. \textbf{EFCore.NamingConventions} is a
  small, clean convention plugin. The \textbf{NetTopologySuite} plugin in the EF Core
  repository shows type mappings and translations. Npgsql's \textbf{NodaTime} plugin shows a
  provider-specific plugin activated through the provider's options builder.
\end{tip}

% ------------------------------------------------------------
\chapter{EF Core's architecture from a plugin's point of view}\label{ch:ef-core-s-architecture-from}

\section{Two containers: the internal service provider}

Every \code{Db\allowbreak{}Context} resolves its infrastructure from an \textbf{internal}
\code{IService\allowbreak{}Provider}. It is separate from the application's DI container and
contains several hundred services: type mapping, model building, the query pipeline, change
tracking, the update pipeline, migrations. A plugin hooks in by registering additional
services in this container or replacing existing ones.

\begin{flowdiagram}
  \flowstep{DbContextOptions}{Provider extension (UseNpgsql) and plugin extensions
  (UseAcmeTenancy)}
  \flowstep{ServiceProviderCache}{Finds a matching internal provider by hash, or builds one}
  \flowstep{ApplyServices()}{Each extension registers its services in the new provider}
  \flowstep{Internal service provider}{Cached and shared by all matching contexts}
  \flowstep{Model (cached)}{Built once per provider: conventions, type mappings, validation}
  \flowstep{Query, SaveChanges, Migrations}{Run with the services of this provider}
\end{flowdiagram}

Two caches are decisive here. The \textbf{service provider} is cached based on the options:
contexts with the same configuration share one provider. The \textbf{model} is cached per
provider (and context type). Both make EF Core fast, but they have direct consequences for
plugins, which chapter~\ref{ch:the-options-extension} covers in detail.

\section{Lifetimes}

Every EF service has a fixed lifetime. \textbf{Singleton} here means: once per internal
service provider, not per application. \textbf{Scoped} means: once per context instance. A
singleton must not hold a scoped service, and a service that needs different data per
context must not be a singleton. Appendix~\ref{app:A} lists the lifetimes of all extension
points.

\section{The dependencies pattern}

EF services do not receive their dependencies individually, but as a parameter object, such
as \code{Migrations\allowbreak{}Sql\allowbreak{}Generator\allowbreak{}Dependencies} or
\code{Provider\allowbreak{}Convention\allowbreak{}Set\allowbreak{}Builder\allowbreak{}Dependencies}.
When you derive from an EF service, you accept this object in the constructor and pass it
on. EF Core adds new dependencies to the object without breaking constructors. For plugins
this is a blessing: derived services stay compatible across minor versions.

\begin{sourcecode}[Migrations/\allowbreak{}Internal/\allowbreak{}Tenancy\allowbreak{}Migrations\allowbreak{}Sql\allowbreak{}Generator.\allowbreak{}cs (\allowbreak{}header)]{csharp}
public class TenancyMigrationsSqlGenerator : NpgsqlMigrationsSqlGenerator
{
    public TenancyMigrationsSqlGenerator(
        MigrationsSqlGeneratorDependencies dependencies,     // parameter object from EF
        INpgsqlSingletonOptions npgsqlSingletonOptions)      // additionally from the provider
        : base(dependencies, npgsqlSingletonOptions)
    {
    }
\end{sourcecode}

\begin{note}[Code style in the sample]
  All classes in the sample use classic constructors with \code{private readonly} fields and
  \code{Argument\allowbreak{}Null\allowbreak{}Exception.\allowbreak{}Throw\allowbreak{}If\allowbreak{}Null},
  no primary constructors. Their parameters are mutable throughout the class body, offer no
  natural place for validation and hide which state is stored. Classic constructors also
  match the style of EF Core and Npgsql, whose classes a plugin derives from. The only
  exception is positional records such as \code{Rls\allowbreak{}Policy}, where the syntax
  generates real \code{init} properties.
\end{note}

\section{Extend or replace?}

There are two fundamentally different ways to extend EF Core. \textbf{Plugin interfaces}
such as \code{IMethod\allowbreak{}Call\allowbreak{}Translator\allowbreak{}Plugin} or
\code{IConvention\allowbreak{}Set\allowbreak{}Plugin} are designed for exactly this. They
can be registered multiple times, EF queries all registrations, and several plugins coexist.
\textbf{Replacing} means swapping a single service such as
\code{IMigrations\allowbreak{}Sql\allowbreak{}Generator} for a derived class. That is
sometimes unavoidable, but only one replacement can win per context.

\begin{important}[Plugin interfaces before replacement]
  Always check first whether a plugin interface exists. Replacing a single service collides
  with every other plugin that replaces the same service and forces you to derive from the
  provider's implementation. That ties the plugin to one provider and its internal API.
  Chapter~\ref{ch:registering-services} shows how to stay composable anyway.
\end{important}

\section{Warm-up: a complete plugin in five pieces}

Before the large sample, here is the smallest plugin that is still complete. It makes
PostgreSQL's \code{reverse(\allowbreak{})} available in LINQ as
\code{EF.\allowbreak{}Functions.\allowbreak{}Reverse(\allowbreak{}.\allowbreak{}.\allowbreak{}.\allowbreak{})}.
Every later chapter elaborates one of its five pieces:

\begin{sourcecode}[tests/\allowbreak{}Acme.\allowbreak{}Entity\allowbreak{}Framework\allowbreak{}Core.\allowbreak{}Tenancy.\allowbreak{}Tests/\allowbreak{}Minimal\allowbreak{}Plugin\allowbreak{}Tests.\allowbreak{}cs (\allowbreak{}1/\allowbreak{}2,\allowbreak{} shortened)]{csharp}
// 1. The function users call in LINQ; it is never executed
public static class ReverseDbFunctionsExtensions
{
    public static string Reverse(this DbFunctions _, string value)
        => throw new InvalidOperationException("Reverse can only be used in LINQ queries.");
}

// 2. The translator turns the method call into a SQL function
//    (chapter "Translating methods and members")
public class ReverseTranslator : IMethodCallTranslator
{
    // Constructor stores the ISqlExpressionFactory in _sql

    public SqlExpression? Translate(
        SqlExpression? instance, MethodInfo method, IReadOnlyList<SqlExpression> arguments,
        IDiagnosticsLogger<DbLoggerCategory.Query> logger)
        => method == ReverseMethod
            ? _sql.Function("reverse", [arguments[1]], nullable: true,
                argumentsPropagateNullability: [true], typeof(string),
                arguments[1].TypeMapping)
            : null;
}

// 3. The plugin hands the translator to EF (chapter "Translating methods and members")
public class ReverseTranslatorPlugin : IMethodCallTranslatorPlugin
{
    public ReverseTranslatorPlugin(ISqlExpressionFactory sql)
    {
        Translators = [new ReverseTranslator(sql)];
    }

    public IEnumerable<IMethodCallTranslator> Translators { get; }
}
\end{sourcecode}

\begin{sourcecode}[tests/\allowbreak{}Acme.\allowbreak{}Entity\allowbreak{}Framework\allowbreak{}Core.\allowbreak{}Tenancy.\allowbreak{}Tests/\allowbreak{}Minimal\allowbreak{}Plugin\allowbreak{}Tests.\allowbreak{}cs (\allowbreak{}2/\allowbreak{}2,\allowbreak{} shortened)]{csharp}
// 4. The options extension registers the plugin in EF's internal provider
//    (chapters "The options extension" and "Registering services")
public class ReverseOptionsExtension : IDbContextOptionsExtension
{
    private DbContextOptionsExtensionInfo? _info;

    public DbContextOptionsExtensionInfo Info => _info ??= new ExtensionInfo(this);

    public void ApplyServices(IServiceCollection services)
        => new EntityFrameworkRelationalServicesBuilder(services)
            .TryAdd<IMethodCallTranslatorPlugin, ReverseTranslatorPlugin>();

    public void Validate(IDbContextOptions options)
    {
    }

    // ExtensionInfo: not a provider, no settings -> hash 0, always the same provider
}

// 5. The Use method activates the plugin
//    (chapter "Public API and integration into the application")
public static DbContextOptionsBuilder UseReverse(this DbContextOptionsBuilder builder)
{
    ((IDbContextOptionsBuilderInfrastructure)builder).AddOrUpdateExtension(
        builder.Options.FindExtension<ReverseOptionsExtension>()
        ?? new ReverseOptionsExtension());
    return builder;
}
\end{sourcecode}

\begin{sourcecode}[Usage]{csharp}
options.UseNpgsql(connectionString).UseReverse();

db.Customers.Where(c => EF.Functions.Reverse(c.Name) == "abc");
// SQL: WHERE reverse(c."Name") = 'abc'
\end{sourcecode}

That is all a working plugin needs. EF finds the translator because the extension registers
it in the internal service provider, and every context configured with
\code{Use\allowbreak{}Reverse(\allowbreak{})} shares that provider. Because the argument
references a column, the function needs no evaluatable expression filter
(chapter~\ref{ch:preventing-client-side-evaluation} explains when it does). The test
\code{Minimal\allowbreak{}Plugin\allowbreak{}Tests} checks the translation and that the same
query fails without the plugin. The tenancy plugin consists of exactly these pieces, only
many more of them.

% ------------------------------------------------------------
\chapter{The sample scenario}\label{ch:the-sample-scenario}

The sample plugin \textbf{Acme.EntityFrameworkCore.Tenancy} makes EF Core multi-tenant with
PostgreSQL. The database itself enforces tenant separation through \textbf{Row Level
Security} (RLS): every table gets a policy that only shows rows of the tenant whose ID is
set in the current database session. Even a faulty query or raw SQL cannot read another
tenant's data.

From the user's point of view it looks like this:

\begin{sourcecode}[User's view]{csharp}
public class Order : ITenantOwned               // marker from the core package
{
    public int Id { get; set; }
    public TenantId TenantId { get; set; }      // strongly typed, automatically uuid
    public decimal Total { get; set; }
}

services.AddAcmeTenancy<HttpTenantProvider>();
services.AddDbContext<ShopContext>((sp, o) => o
    .UseNpgsql(connectionString)
    .UseAcmeTenancy(sp, tenancy => tenancy.ApplyQueryFilter()));

// Querying and saving without any tenant logic in application code
db.Orders.Add(new Order { Total = 42 });         // the plugin sets TenantId
var mine = await db.Orders.ToListAsync();        // RLS filters in the database
\end{sourcecode}

To make this work, the plugin uses every extension point covered in this book:

\begin{datatable}{H{1.20}H{1.50}H{0.30}}
  \textbf{Requirement} & \textbf{Implementation} & \textbf{Ch.} \\
  \tablehead
  Configuration via \code{Use\allowbreak{}Acme\allowbreak{}Tenancy(\allowbreak{})} & Options
  extension, settings record, options builder & \ref{ch:the-options-extension}--\ref{ch:public-api-and-integration-into} \\
  \tablerule
  \code{Tenant\allowbreak{}Id} as a column type without configuration & Type mapping plugin,
  value converter, JSON reader/writer & \ref{ch:type-mappings-and-value-converters} \\
  \tablerule
  \code{EF.\allowbreak{}Functions.\allowbreak{}Current\allowbreak{}Tenant(\allowbreak{})} in
  LINQ & Method translator, evaluatable expression filter & \ref{ch:translating-methods-and-members}--\ref{ch:preventing-client-side-evaluation} \\
  \tablerule
  \code{tenant\allowbreak{}Id.\allowbreak{}Is\allowbreak{}Empty} in LINQ & Member translator
  & \ref{ch:translating-methods-and-members} \\
  \tablerule
  RLS for all \code{ITenant\allowbreak{}Owned} entities & Convention with index and optional
  query filter & \ref{ch:conventions} \\
  \tablerule
  Fluent API, attribute \code{[No\allowbreak{}Tenant\allowbreak{}Isolation]} & Annotations,
  attribute convention, ConfigurationSource & \ref{ch:annotations-and-fluent-api}--\ref{ch:conventions} \\
  \tablerule
  Detect misconfiguration early & Validation convention, custom warnings & \ref{ch:model-validation}, \ref{ch:diagnostics-and-logging} \\
  \tablerule
  Set TenantId automatically & SaveChanges interceptor & \ref{ch:interceptors} \\
  \tablerule
  Tenant in the DB session & Connection interceptor & \ref{ch:interceptors} \\
  \tablerule
  \code{ENABLE ROW LEVEL SECURITY}, \code{CREATE POLICY} & Annotation provider, migrations
  SQL generator & \ref{ch:how-migrations-work}--\ref{ch:declarative-sql-from-annotations} \\
  \tablerule
  Roles and grants in migrations & Custom migration operations & \ref{ch:custom-migration-operations} \\
  \tablerule
  Readable snapshots & Annotation code generator, design-time services & \ref{ch:design-time-services} \\
  \tablerule
  Compiled models and AOT & \code{Default} instances, source generators & \ref{ch:compiled-models-trimming-and-nativeaot} \\
  \tablerule
  \code{EF.\allowbreak{}Functions.\allowbreak{}Is\allowbreak{}Not\allowbreak{}Distinct\allowbreak{}From},
  \code{EF.\allowbreak{}Functions.\allowbreak{}Every} & Custom SQL node, aggregate
  translator & \ref{ch:custom-sql-expression-nodes}--\ref{ch:aggregate-translators} \\
  \tablerule
  Tenant property with any name &
  \code{Is\allowbreak{}Tenant\allowbreak{}Property(\allowbreak{})}, runtime annotation & \ref{ch:property-level-fluent-api-and}
  \\
  \tablerule
  Scaffold RLS from an existing database & Custom database model factory & \ref{ch:scaffolding-in-practice} \\
  \tablerule
  SQL Server & Second provider package on the same core & \ref{ch:what-belongs-in-the-core}--\ref{ch:a-second-provider-sql-server} \\
\end{datatable}

\section{Package architecture}

The plugin consists of three NuGet packages. The \textbf{core package} contains everything
that is independent of the database and that assemblies with entity classes may reference
without pulling in a provider. The \textbf{PostgreSQL package} is the complete
implementation this book uses to explain everything. The \textbf{SQL Server package} uses
the same core and shows in chapter~\ref{ch:a-second-provider-sql-server} what a second
provider has to implement itself.

\begin{sourcecode}[Structure of the three packages]{text}
Acme.EntityFrameworkCore.Tenancy               Core (EF Core relational, no provider)
├─ TenantId, TenantIdConverter, TenantIdJsonReaderWriter, TenancyException
├─ ITenantOwned, ITenantProvider, NoTenantIsolationAttribute
├─ EF.Functions.CurrentTenant(), TenancyDbFunctionsAttribute
├─ Metadata/         runtime annotation, IsTenantProperty(), candidate lookup
├─ Conventions/      value generation for TenantId keys
├─ Query/            evaluatable expression filter (for all packages)
├─ Interceptors/     SaveChanges interceptor, abstract session interceptor
├─ Diagnostics/      event IDs, logging definitions, logger extensions
└─ DependencyInjection/  AddAcmeTenancy<T>()

Acme.EntityFrameworkCore.Tenancy.PostgreSQL    complete implementation
├─ Infrastructure/   options extension, builder, singleton options
├─ Storage/          TenantIdTypeMapping, type mapping plugin
├─ Query/            translators, custom SQL node, aggregate, SQL generator
├─ Metadata/         RLS annotations, RlsPolicy, annotation provider
├─ Builders/         Fluent API (HasRowLevelSecurity, HasPolicy, ...)
├─ Conventions/      isolation, attribute, validation, convention set plugin
├─ Migrations/       SQL generator, custom operations
├─ Design/           code generators, design-time services
├─ Scaffolding/      read RLS and policies from the database
└─ Runtime/          TenantConnectionInterceptor

Acme.EntityFrameworkCore.Tenancy.SqlServer     second provider
                                               (chapter "A second provider: SQL Server")
└─ Infrastructure/, Storage/, Query/, Metadata/, Conventions/, Migrations/, Runtime/
\end{sourcecode}

\begin{tip}
  The folder structure follows EF Core's own (Infrastructure, Storage, Query, Metadata,
  Migrations, Design). Anyone who knows EF Core finds their way immediately, and when
  looking things up in the EF source you know where the counterpart lives.
\end{tip}
